%% chapters/13-eigenvectors.tex
\chapter{Eigenvectors: When They Converge and When They Don't}
\label{ch:eigvecs}

\whereWeAre{%
Convergence of eigenvalues does not imply convergence of eigenvectors. When the limiting eigenvalue is simple, the spectral projector is continuous and a sign convention pins down a converging eigenvector. When eigenvalues coincide, the eigenspace is multi-dimensional and the eigenvectors can rotate freely without changing the matrix limit.%
}

\PLACEHOLDER{Sources: April 17 Gustavo notes (rotation counterexample). Existing main.tex April 17 continuation section.}

\section{Spectral projectors and the continuity argument}
\PLACEHOLDER{Riesz contour formula; isolated eigenvalues give continuous projectors; sign convention.}

\section{When eigenvalues coincide}
\PLACEHOLDER{Eigenspace = whole plane; any basis is an eigenbasis; no canonical limit.}

\section{The rotation counterexample}
\PLACEHOLDER{$S_n = O_n \diag(\lambda_1^{(n)}, \lambda_2^{(n)}) O_n^T$ with $\theta_n = n\pi/2$; $S_n \to I_2$ in operator norm but the eigenvectors rotate.}

\section{Brief warm-ups}
\PLACEHOLDER{3--5 short questions.}

\chapterRecap{%
\begin{itemize}
  \item \PLACEHOLDER{Eigenvector convergence requires eigenvalue isolation.}
  \item \PLACEHOLDER{The rotation counterexample makes the failure tangible.}
\end{itemize}%
}

\section*{Exercises}
\PLACEHOLDER{8--15 longer exercises.}
